\section{Weekly Summary}

\begin{tcolorbox}[colback=yellow!10,colframe=black!50,title=AI-Generated Content]
\noindent The summary below was written by Claude (Anthropic) from that week's regression outputs and series descriptions. It has not been reviewed or fact-checked by a human, and should be read as an automated, informal recap -- not analysis.
\end{tcolorbox}

Welcome back to another thrilling week of forcing unrelated economic data to hold hands and smile for the camera. As always, the premise is simple: we grab two series that have no business being in the same room together, run a line through them, and then act shocked at whatever falls out. This week's lineup included a World Pandemic Uncertainty Index for Myanmar, a discontinued San Francisco tech barometer, textile mill inventories, auto production seasonal factors, and -- my personal favorite -- the dollars spent by certain racial groups on food during out-of-town trips. Science, baby.

Monday kicked things off with a bang: rest-of-the-world debt securities versus per-capita housing-and-utilities spending in the Plains region posted a jaw-dropping R-squared of 0.96. Ninety-six percent! Call the Nobel committee! Except, oh no, when we detrended both series and stripped out the fact that they're both just marching upward over time like two escalators in the same mall, that glorious 0.96 deflated into a flabby 0.15 that barely clears the significance bar by holding onto the railing. Classic trend mirage. The two series weren't dancing together; they were just both getting taller. Thursday served up the same sad story -- federal debt held by the Fed versus textile mill inventories looked gorgeous at 0.65 raw, then collapsed to a thoroughly unconvincing 0.02 once the shared drift was removed. Shocking that "two numbers that grew over the decades" grew over the decades.

The rest of the week was a parade of relationships so weak they'd lose an arm-wrestling match to a damp napkin. Household appliance import prices versus syndicated loan percentages? Nothing, raw or detrended. Myanmar's pandemic uncertainty versus San Francisco tech? A big fat nope. Aircraft manufacturing wages versus Channel Islands bond valuation changes looked briefly promising raw at 0.23, then detrended itself into complete irrelevance, even flipping the direction of the relationship for good measure, which is nature's way of saying "there was never anything here." And the grand finale, the real user cost of some money-supply index versus auto production seasonal factors, achieved an R-squared of 0.0002 both ways -- a result so flat you could use it as a spirit level.

So, to recap the recap: zero relationships this week survived detrending with their dignity intact, which is honestly the house special around here. Let me remind you, with all the gravitas I can muster, that these series are paired at random and none of this means anything -- no causation, no theory, no actual link, just two time series that got unlucky enough to be drawn from the same hat on the same day. Spurious correlation is not a bug of this project; it is the entire personality. See you next week for more rigorous nonsense.
